% TOP_PROBLEM: {"id": 2305005, "title": "Research Problems in Function Theory — Problem 5.5", "category_id": 8, "category": {"id": 8, "name": "analysis", "display_name": "Analysis", "description": "Limits, continuity, calculus, and function theory.", "slug": "analysis", "order_index": 8, "created_at": "2026-07-31T15:26:25.671Z"}, "status": "open"}
% FIRST_POSTED: null
% ENTRY_KIND: statement_only
% Imported from: batch14.tex; UnsolvedMath ID 2305005
% Compile twice: lualatex 2305005.tex
\documentclass[11pt]{article}
\usepackage{fontspec}
\setmainfont{Latin Modern Roman}
\usepackage{amsmath,amssymb,mathtools,unicode-math}
\setmathfont{Latin Modern Math}
\usepackage{xurl,hyperref}
\hypersetup{hidelinks,pdftitle={UnsolvedMath Problem 2305005}}
\newcommand{\sref}[2]{\hyperlink{source:#1:#2}{[S#2]}}
\newcommand{\cataloguescope}[1]{\paragraph{Mathematical question.}#1}
\begin{document}
% UnsolvedMath ID 2305005; source catalogue code AMR-022-5005
\setcounter{section}{2305004}
\section{Research Problems in Function Theory — Problem 5.5}\label{2305005}
{\small\sffamily UnsolvedMath ID 2305005 · Analysis}
\subsection{Definitions and mathematical statement}

\paragraph{Shared notation.}$\mathbb N=\{1,2,\ldots\}$, and $\mathbb Z,\mathbb Q,\mathbb R,\mathbb C$ denote the integers, rationals, reals and complex numbers. Any other conventions are specified locally.


Write $\mathbb D=\{z\in\mathbb C:|z|<1\}$ and $\mathbb T=\partial\mathbb D$. All Taylor coefficients below are taken at $0$. For holomorphic $f(z)=\sum_{n=0}^\infty a_nz^n$ on $\mathbb D$, put
\[d_f(R)=\sup\{\rho>0:B(w,\rho)\subset f(\mathbb D)\text{ for some }|w|=R\}\quad(R>0),\]
where $B(w,\rho)=\{v:|v-w|<\rho\}$ and the supremum of the empty set is $0$. Thus this measures the largest full image disc whose center has modulus $R$. The assertion $a_n=O(1)$ means $\sup_{n\ge1}|a_n|<\infty$ for the particular function.
\paragraph{Statement recorded in the supplied source.}
Describe conditions on an image domain $D\subset\mathbb C$ ensuring that every holomorphic $f:\mathbb D\to D$ has $a_n=O(1)$. In particular, does boundedness of the coefficients still follow when $d_f(R)$ is allowed to tend to infinity sufficiently slowly? Formulate the allowable growth rates: determine whether there is a function $h:[1,\infty)\to(0,\infty)$, with $h(R)\to\infty$, for which $d_f(R)\le h(R)$ for all sufficiently large $R$ forces $a_n=O(1)$ for every such $f$. The bound for the coefficients may depend on $f$. \sref{2305005}{2}
\subsection{Short English statement}
Which restrictions on the image of a disc function force bounded Taylor coefficients, and can the size of full image discs grow slowly without losing that conclusion?
\subsection{Sources}
\begin{description}
\item[\hypertarget{source:2305005:1}{S1}] ULAM.AI and the UnsolvedMath Contributors, UnsolvedMath: A Curated Collection of Open Mathematics Problems, record 2305005 (AMR-022-5005). CC BY 4.0. \url{https://huggingface.co/datasets/ulamai/UnsolvedMath}.
\item[\hypertarget{source:2305005:2}{S2}] W. K. Hayman and E. F. Lingham, Research Problems in Function Theory (2018), arXiv:1809.07200, Chapter 5, Problem 5.5, its update and the preceding shared notation. \url{https://arxiv.org/abs/1809.07200}.
\end{description}
\label{end:2305005}
\end{document}
